\chapter{Sequential Monte Carlo Methods}

Sequential Monte Carlo (SMC) methods provide a flexible framework for approximating probability distributions that evolve sequentially over time. 
They are particularly useful for inference in state-space models, where the latent states are typically not directly observed and their posterior distributions are analytically intractable. 
SMC methods approximate these distributions using a set of weighted samples, commonly referred to as particles, and update this approximation as new observations become available.

This chapter first introduces the foundations of SMC, beginning with Monte Carlo methods and sequential importance sampling. 
The problem of weight degeneracy is then addressed through resampling, leading to sequential importance sampling with resampling and particle filtering methods. 
The bootstrap particle filter and auxiliary particle filter are then presented in more detail.

The chapter then considers particle based methods for Bayesian parameter inference. 
Particle marginal Metropolis-Hastings and Particle Gibbs are introduced as particle Markov chain Monte Carlo methods. 
Finally, SMC$^2$ is presented as an extension that combines particle filtering with sequential inference for static parameters and latent states.

\section{Monte Carlo method}
Monte Carlo methods approximate quantities involving a probability distribution by using a finite collection of random samples. The basic idea is simple: rather than evaluating an integral directly, we replace it by an average over samples drawn from the distribution of interest.

Until further notice, let $X$ be a random variable with probability density function $p(x)$, and let $g$ be a measurable function such that $E_p[|g(X)|] < \infty$. The quantity of interest is
\begin{equation}
    I = E_p[g(X)] = \int g(x)p(x)\,\mathrm{d}x.
    \label{eq:target_expectation}
\end{equation}

\begin{definition}
    Let $X_1,\ldots,X_N \overset{\mathrm{iid}}{\sim} p(x)$. The \emph{Monte Carlo estimator} of $I$ is defined as
    \begin{equation}
        \hat{I}_N = \frac{1}{N}\sum_{i=1}^{N} g(X_i).
        \label{eq:monte_carlo_estimator}
    \end{equation}
\end{definition}

Such an estimator has several desirable properties, including unbiasedness, consistency, and asymptotic normality.

\begin{proposition}
    The Monte Carlo estimator $\hat{I}_N$ is unbiased.
\end{proposition}
\begin{proof}
    The proof follows from the linearity of expectation and the independence of the samples.
    \begin{align}
        E_p[\hat{I}_N] &= E_p\left[\frac{1}{N}\sum_{i=1}^{N} g(X_i)\right] \\
        &= \frac{1}{N}\sum_{i=1}^N E_p[g(X_i)] \\
        &= \frac{1}{N}\sum_{i=1}^N I = I.
    \end{align}
\end{proof}

\begin{proposition}
    The Monte Carlo estimator $\hat{I}_N$ is consistent.
\end{proposition}
\begin{proof}
    Because $E_p[g(X)] = I$ is finite, the strong law of large numbers applies, giving
    \begin{equation}
        \hat{I}_N \xrightarrow[N\to\infty]{a.s.} I.
    \end{equation}
    Almost sure convergence implies convergence in probability, so the estimator is consistent.
\end{proof}

\begin{proposition}
    If $g(X)$ has finite variance, then the Monte Carlo estimator $\hat{I}_N$ is asymptotically normal.
\end{proposition}
\begin{proof}
    If $g(X)$ has finite variance $\sigma^2$, then the central limit theorem applies, giving
    \begin{equation}
        \sqrt{N}\left(\hat{I}_N-I\right)
        \xrightarrow[N\to\infty]{D}
        \mathcal{N}(0,\sigma^2),
    \end{equation}
    which is the definition of asymptotic normality.
\end{proof}

\subsection{Empirical distribution}
\label{sec:empirical_distribution}

The Monte Carlo estimator can be interpreted as integration with respect to an empirical approximation constructed from the samples.

\begin{definition}
Given samples $X_1,\ldots,X_N \overset{\mathrm{iid}}{\sim} p(x)$, the \emph{empirical distribution} associated with these samples is defined as
\begin{equation}
    \hat{p}_N(x)
    = \frac{1}{N}\sum_{i=1}^{N}\delta(x-X_i),
    \label{eq:empirical_distribution}
\end{equation}
where $\delta(\cdot)$ denotes the Dirac delta function.
\end{definition}

The empirical distribution assigns equal weight $1/N$ to each sampled value. It therefore provides a discrete approximation to the distribution represented by $p(x)$. The expression in~\eqref{eq:empirical_distribution} should be understood through integration rather than as an ordinary probability density with respect to Lebesgue measure.

\begin{proposition}
For any measurable function $g$ with finite expectation,
\begin{equation}
    \hat{I}_N = \int g(x)\hat{p}_N(x)\,\mathrm{d}x
\end{equation}
\end{proposition}

\begin{proof}
    \begin{align}
        \int g(x)\hat{p}_N(x)\,\mathrm{d}x &= \int g(x) \frac{1}{N}\sum_{i=1}^N \delta(x-X_i)\,\mathrm{d}x \\
        &= \frac{1}{N}\sum_{i=1}^N \int g(x)\delta(x-X_i)\,\mathrm{d}x \\
        &= \frac{1}{N}\sum_{i=1}^N g(X_i) = \hat{I}_N,
    \end{align}
    where the defining property of the Dirac delta function has been used.
\end{proof}

It is important to distinguish the empirical distribution $\hat{p}_N$ from the density $p$. Even when $p$ is continuous, $\hat{p}_N$ consists of point masses for every finite $N$ and is therefore not a pointwise approximation of $p(x)$. The approximation is instead meaningful through quantities such as integrals and expectations.

\section{Importance sampling}
\label{sec:importance_sampling}
The standard Monte Carlo method requires samples from the target distribution $p$. In many problems, however, sampling directly from $p$ is difficult or even impossible.
Importance sampling addresses this problem by introducing a second distribution from which sampling is easier.

\begin{definition}
Let $q(x)$ be a proposal distribution satisfying $q(x)>0$ whenever $p(x)>0$. 
The \emph{importance weight} associated with a point $x$ is defined as
\begin{equation}
    w(x) = \frac{p(x)}{q(x)}.
    \label{eq:importance_weight}
\end{equation}
\end{definition}

The expectation of $g$ under $p$ can then be rewritten as
\begin{align}
    I = E_p[g(X)] &= \int g(x)p(x)\,\mathrm{d}x \\
    &= \int g(x)\frac{p(x)}{q(x)}q(x)\,\mathrm{d}x \\
    &= E_q\left[g(X)w(X)\right].
\end{align}
Thus, although samples are generated from $q$, the importance weights correct for the difference between $q$ and $p$.

\begin{definition}
Let $X_1,\ldots,X_N \overset{\mathrm{iid}}{\sim} q$ and $w(X_i)$ corresponding importance weights.
The \emph{importance sampling estimator} of $I=E_p[g(X)]$ is defined by
\begin{equation}
    \hat{I}_N^{\mathrm{IS}} = \frac{1}{N} \sum_{i=1}^N g(X_i)w(X_i).
    \label{eq:importance_sampling_estimator}
\end{equation}
\end{definition}
\begin{proposition}
    \label{prop:importance_sampling_unbiased}
    The importance sampling estimator $\hat{I}_N^{\mathrm{IS}}$ is unbiased.
\end{proposition}
\begin{proof}
    \begin{align}
        E_q[\hat{I}_N^{\mathrm{IS}}] &= \frac{1}{N}\sum_{i=1}^N E_q[g(X_i)w(X_i)] \\
        &= E_q\left[g(X)\frac{p(X)}{q(X)}\right] \\
        &= \int g(x)\frac{p(x)}{q(x)}q(x)\,\mathrm{d}x = I.
    \end{align}
\end{proof}

\begin{proposition}
    \label{prop:importance_sampling_consistent}
    The importance sampling estimator $\hat{I}_N^{\mathrm{IS}}$ is consistent.
\end{proposition}
\begin{proof}
    Because $E_q[g(X)w(X)] = I$ is finite, the strong law of large numbers gives
    \begin{equation}
        \hat{I}_N^{\mathrm{IS}} \xrightarrow[N\to\infty]{a.s.} I,
    \end{equation}
    which means that the estimator is consistent.
\end{proof}

Although the importance sampling estimator is unbiased and consistent for any proposal distribution satisfying the support condition, its variance depends strongly on the choice of $q$.

\subsection{(Weighted) empirical distribution}
\label{sec:weighted_empirical_distribution}

There are two equivalent ways to interpret the importance sampling estimator.
The first one follows the same logic as in \ref{sec:empirical_distribution}:
\begin{align}
    \hat{I}_N^{\mathrm{IS}} &= \frac{1}{N}\sum_{i=1}^{N}w(X_i)g(X_i) = \frac{1}{N}\sum_{i=1}^{N} \int w(x)g(x) \delta(x-X_i)\,\mathrm{d}x \\
    &=\int w(x)g(x) \frac{1}{N}\sum_{i=1}^{N}\delta(x-X_i)\,\mathrm{d}x \\
    &=\int w(x)g(x) \hat{q}_N(x)\,\mathrm{d}x,
\end{align}
where samples $X_1,\ldots,X_N\overset{\mathrm{iid}}{\sim}q$ construct an empirical distribution $\hat{q}_N(x)$ and we are simply computing the expectation of $w(X)g(X)$ with respect to it.

There is another interpretation that is particularly useful. Instead of viewing the weights as part of the integrand, we can absorb them into the empirical distribution itself.
\begin{align}
    \hat{I}_N^{\mathrm{IS}} &= \frac{1}{N}\sum_{i=1}^{N}w(X_i)g(X_i) = \frac{1}{N}\sum_{i=1}^{N} \int w(x)g(x) \delta(x-X_i)\,\mathrm{d}x \\
    &=\int g(x) \frac{1}{N}\sum_{i=1}^{N}w(x)\delta(x-X_i)\,\mathrm{d}x \\
    &=\int g(x) \hat{p}_N^{\mathrm{IS}}(x)\,\mathrm{d}x,
\end{align}
leading us to next definition.
\begin{definition}
Let $X_1,\ldots,X_N \overset{\mathrm{iid}}{\sim} q$ and $w(X_i)$ corresponding importance weights. 
The \emph{weighted empirical distribution} is defined as
\begin{equation}
    \hat{p}_N^{\mathrm{IS}}(x) = \frac{1}{N} \sum_{i=1}^N w(X_i)\delta(x-X_i).
    \label{eq:weighted_empirical_distribution}
\end{equation}
\end{definition}

This provides an alternative interpretation of importance sampling:
although the samples are drawn from the proposal distribution $q$, the
importance weights allow the samples to represent the target distribution
$p$. Unlike ordinary Monte Carlo, where every sample contributes the same
mass $1/N$, importance sampling assigns the sample $X_i$ a mass
proportional to $w(X_i)$.

The weighted empirical distribution in
\eqref{eq:weighted_empirical_distribution}, however, is not generally a
probability distribution for a finite sample. Its total mass is
\begin{equation}
    \int \hat{p}_N^{\mathrm{IS}}(x)\,\mathrm{d}x = \frac{1}{N} \sum_{i=1}^N w(X_i),
\end{equation}
which is a random quantity and is not necessarily equal to one.
Nevertheless, the total mass has expectation equal to one. Since $X\sim q$,
\begin{align}
    E_q[w(X)] &= \int \frac{p(x)}{q(x)}q(x)\,\mathrm{d}x = \int p(x)\,\mathrm{d}x = 1, \\
    E_q\left[\frac{1}{N}\sum_{i=1}^N w(X_i)\right] &= \frac{1}{N}\sum_{i=1}^N E_q[w(X_i)] = 1.
\end{align}
Additionally, the strong law of large numbers gives
\begin{equation}
    \frac{1}{N}\sum_{i=1}^N w(X_i) \xrightarrow[N\to\infty]{a.s.} E_q[w(X)] = 1.
\end{equation}
Thus, although the weighted empirical distribution does not necessarily
have unit mass for a finite sample, its total mass converges almost surely
to one as the number of samples increases.

\section{Self-normalized importance sampling}
\label{sec:self_normalized_importance_sampling}
The previous construction suggests a natural modification: instead of allowing the weights to have a random total mass, normalize them as
\begin{equation}
    \tilde{w}_i = \frac{w(X_i)}{\sum_{j=1}^{N} w(X_j)},
\end{equation}
which leads us to the next definition.
\begin{definition}
    Let $X_1,\ldots,X_N \overset{\mathrm{iid}}{\sim} q$ and $\tilde{w}(X_i)$ corresponding normalized importance weights.
    The \emph{self-normalized weighted empirical distribution} is defined as
    \begin{equation}
        \hat{p}_N^{\mathrm{SNIS}}(x) = \sum_{i=1}^{N} \tilde{w}_i \delta(x-X_i)
    \end{equation}
\end{definition}
This distribution clearly has a total mass of one.
\begin{definition}
    Let $X_1,\ldots,X_N \overset{\mathrm{iid}}{\sim} q$ and $w(X_i)$ corresponding importance weights.
    The \emph{self-normalized imporance sampling estimator} of $I=E_p[g(X)]$ is defined as
    \begin{equation}
        \hat{I}_N^{\mathrm{SNIS}} = \frac{\sum_{i=1}^{N} g(X_i)w(X_i)}{\sum_{i=1}^{N}w(X_i)}.
    \end{equation}
\end{definition}

\begin{proposition}
    \label{prop:self_normalized_importance_sampling_unbiased}
    For any measurable function $g$ with finite expectation,
    \begin{equation}
        \hat{I}_N^{\mathrm{SNIS}} = \int g(x) \hat{p}_N^{\mathrm{SNIS}}\, \mathrm{d}x.
    \end{equation}
\end{proposition}
\begin{proof}
    \begin{align}
        \int g(x) \hat{p}_N^{\mathrm{SNIS}}\, \mathrm{d}x &= \int g(x) \sum_{i=1}^{N} \tilde{w}_i \delta(x-X_i)\,\mathrm{d}x \\
        &= \sum_{i=1}^{N} \tilde{w}_i \int g(x) \delta(x-X_i)\,\mathrm{d}x \\
        &= \sum_{i=1}^{N} \tilde{w}_i g(X_i) \\
        &= \frac{\sum_{i=1}^{N} g(X_i)w(X_i)}{\sum_{i=1}^{N}w(X_i)}.
    \end{align}
\end{proof}

There is a fundamental difference between $\hat{I}_N^{\mathrm{IS}}$ and $\hat{I}_N^{\mathrm{SNIS}}$. The ordinary importance sampling estimator is unbiased, whereas the self-normalized estimator is generally biased for finite $N$.
\begin{proposition}
    The self-normalized importance sampling estimator is biased.
\end{proposition}
\begin{proof}
    Because the expected value of a fraction is not equal to the fraction of expected values we have
    \begin{align}
        E_q\left[\frac{\sum_{i=1}^{N}g(X_i)w(X_i)}{\sum_{i=1}^{N}w(X_i)}\right] &\neq \frac{E_q\left[\frac{1}{N} \sum_{i=1}^{N}g(X_i)w(X_i)\right]}{E_q\left[\frac{1}{N} \sum_{i=1}^{N}w(X_i)\right]} \\
        E_q\left[\hat{I}_N^{\mathrm{SNIS}}\right] &\neq \frac{E_q\left[\hat{I}_N^{\mathrm{IS}}\right]}{1} = I.
    \end{align}
\end{proof}

\begin{proposition}
    The self-normalized importance sampling estimator is consistent.
\end{proposition}
\begin{proof}
    Under suitable regularity conditions, the strong law of large numbers gives
    \begin{align}
        \frac{1}{N}\sum_{i=1}^{N} g(X_i)w(X_i) &\xrightarrow[N\to\infty]{a.s.} I,\\
        \frac{1}{N}\sum_{i=1}^{N} w(X_i) &\xrightarrow[N\to\infty]{a.s.} 1.
    \end{align}
    Consequently,
    \begin{equation}
        \hat{I}_N^{\mathrm{SNIS}} = \frac{\frac{1}{N}\sum_{i=1}^{N} g(X_i)w(X_i)}{\frac{1}{N}\sum_{i=1}^{N} w(X_i)} \xrightarrow[N\to\infty]{a.s.} I.
    \end{equation}
\end{proof}

At first it seems pointless to use the self-normalized estimator, just to assure that the $\hat{p}_N^{\mathrm{SNIS}}$ is a proper probability distribution, as we lose the unbiasedness of the estimator. 
However, in many applications, especially in Bayesian inference, the target density $p(x)$ is known only up to a proportionality constant. Suppose that our target has form
\begin{equation}
    p(x) = \frac{\tilde{p}(x)}{Z},
\end{equation}
where $\tilde{p}(x)$ is known and $Z$ is an intractable normalization constant. The importance weights are then
\begin{equation}
    w(x) = \frac{p(x)}{q(x)} = \frac{\tilde{p}(x)}{Zq(x)},
\end{equation}
which makes it impossible to compute the ordinary importance sampling estimator $\hat{I}_N^{\mathrm{IS}}$. However, the normalized weights are still computable:
\begin{equation}
    \tilde{w}_i = \frac{w(X_i)}{\sum_{j=1}^N w(X_j)} = \frac{\frac{\tilde{p}(X_i)}{Zq(X_i)}}{\sum_{j=1}^N \frac{\tilde{p}(X_j)}{Zq(X_j)}} = \frac{\frac{\tilde{p}(X_i)}{q(X_i)}}{\sum_{j=1}^N \frac{\tilde{p}(X_j)}{q(X_j)}}.
\end{equation}
Therefore, the self-normalized importance sampling estimator $\hat{I}_N^{\mathrm{SNIS}}$ can be used even when the target density is known only up to a constant of proportionality at the cost of introducing bias.

\section{Sequential importance sampling}
Sequential importance sampling (SIS) applies importance sampling to a sequence of random variables by exploiting the sequential structure of the problem.
Because the main focus of this thesis are state-space models, in particular the MSSV model, we will consider the sequential importance sampling in the context of state-space models.

Let's assume the model parameters $\theta$ are known for the moment. Given the observations $y_{1:T}$, the objective is to approximate the posterior distribution of the latent states,
\begin{equation}
p(x_{0:T} \mid y_{1:T},\theta) = \frac{p(x_{0:T},y_{1:T}\mid\theta)}{p(y_{1:T}\mid\theta)},
\end{equation}
where
\begin{equation}
p(y_{1:T}\mid\theta) = \int p(x_{0:T},y_{1:T}\mid\theta)\,\mathrm{d}x_{0:T}
\end{equation}
is the \emph{marginal likelihood}. The marginal likelihood is generally intractable for nonlinear and non-Gaussian state-space models.

For notational simplicity, the notation for dependence on $\theta$ will be omitted. As derived in Section~\ref{sec:problem_formulation}, the joint density is
\begin{equation}
p(x_{0:T},y_{1:T}) = \chi(x_0) \prod_{t=1}^T f(y_t\mid x_t) p(x_t\mid x_{t-1}).
\end{equation}

Let's assume that the proposal distribution can also be factorized sequentially as
\begin{equation}
q(x_{0:T}\mid y_{1:T}) = q_0(x_0) \prod_{t=1}^T q_t(x_t\mid x_{0:t-1},y_{1:t}),
\end{equation}
from which we can sample entire trajectories of the latent states, conditioned on the currently available observations, in a sequential manner.

The importance weight can therefore be computed recursively as
\begin{align}
    w_T &= \frac{p(x_{0:T},y_{1:T})}{q(x_{0:T}\mid y_{1:T})} = \frac{\chi(x_0) \prod_{t=1}^T f(y_t\mid x_t) p(x_t\mid x_{t-1})}{q_0(x_0) \prod_{t=1}^T q_t(x_t\mid x_{0:t-1},y_{1:t})} \\
    &= w_{T-1} \cdot \frac{f(y_T\mid x_T) p(x_T\mid x_{T-1})}{q_T(x_T\mid x_{0:T-1},y_{1:T})},
\end{align}
with and $w_0 = \frac{\chi(x_0)}{q_0(x_0)}$.

Let $X_{0:T}^{(i)} \overset{\text{iid}}{\sim} q$ be $N$ independent trajectories sampled from the proposal distribution with $w_T^{(i)}$ their corresponding importance weights.
The marginal likelihood can be estimated using the importance weights of sampled trajectories as
\begin{equation}
    \hat{p}(y_{1:T}) = \frac{1}{N} \sum_{i=1}^{N} w_T^{(i)}.
    \label{eq:marginal_likelihood_estimator}
\end{equation}
\begin{proposition}
    The marginal likelihood estimator $\hat{p}(y_{1:T})$ is unbiased and consistent.
    \label{prop:marginal_likelihood_estimator}
\end{proposition}
\begin{proof}
    Unbiasedness follows from
    \begin{align}
        E_q[\hat{p}(y_{1:T})] &= E_q\left[\frac{1}{N} \sum_{i=1}^{N} w_T^{(i)}\right] = \frac{1}{N} \sum_{i=1}^{N} E_q[w_T^{(i)}] \\
        &= E_q[w_T] = \int w_T q(x_{0:T}\mid y_{1:T})\,\mathrm{d}x_{0:T} \\
        &= \int p(x_{0:T},y_{1:T})\,\mathrm{d}x_{0:T} = p(y_{1:T}),
    \end{align}
    where the linearity of expectation and the independence of the samples have been used.

    Alternatively, $\hat{p}(y_{1:T})$ is just a concrete example of (ordinary) importance sampling estimator $\hat{I}_N^{\mathrm{IS}}$, which we already proved to be unbiased in Proposition~\ref{prop:importance_sampling_unbiased}. Therefore, the consistency of $\hat{p}(y_{1:T})$ also follows from Proposition~\ref{prop:importance_sampling_consistent}.
\end{proof}

Now we can move on to the smoothing distribution, specifically the expectation of a function $g$ under the smoothing distribution, which is
\begin{align*}
    I = E\left[g(x_{0:T}) \mid Y_{1:T} \right] &= \int g(x_{0:T}) p(x_{0:T}\mid y_{1:T})\,\mathrm{d}x_{0:T} \\
    &= \frac{1}{p(y_{1:T})} \int g(x_{0:T}) \frac{p(x_{0:T},y_{1:T})}{q(x_{0:T}\mid y_{1:T})} q(x_{0:T}\mid y_{1:T})\,\mathrm{d}x_{0:T} \\
    &= \frac{1}{p(y_{1:T})} \int g(x_{0:T}) w_T(x_{0:T})q(x_{0:T}\mid y_{1:T}) \,\mathrm{d}x_{0:T}.
\end{align*}
Natural approximations of $I$ can be obtained by replacing the integral with a Monte Carlo average over the sampled trajectories, hence obtaining the estimator
\begin{equation}
    \hat{I}_N = \frac{\frac{1}{N} \sum_{i=1}^{N} g(X_{0:T}^{(i)}) w_T^{(i)}}{p(y_{1:T})} = \frac{\sum_{i=1}^{N} g(X_{0:T}^{(i)}) w_T^{(i)}}{\sum_{i=1}^{N} w_T^{(i)}} = \sum_{i=1}^{N} g(X_{0:T}^{(i)}) \tilde{w}_T^{(i)},
    \label{eq:smoothing_estimator}
\end{equation}
where $\tilde{w}_T^{(i)}$ are normalized weights. Clearly, this is a self-normalized importance sampling estimator, hence it is biased but consistent. Since
\begin{align}
    \hat{I}_N &= \sum_{i=1}^{N} g(X_{0:T}^{(i)}) \tilde{w}_T^{(i)} = \sum_{i=1}^{N} \tilde{w}_T^{(i)} \int g(x_{0:T}) \delta(x_{0:T}-X_{0:T}^{(i)})\,\mathrm{d}x_{0:T} \\
    &= \int g(x_{0:T}) \sum_{i=1}^{N} \tilde{w}_T^{(i)} \delta(x_{0:T}-X_{0:T}^{(i)})\,\mathrm{d}x_{0:T} = \int g(x_{0:T}) \hat{p}_N(x_{0:T}\mid y_{1:T})\,\mathrm{d}x_{0:T},
\end{align}
the empirical distribution of the smoothing distribution is
\begin{equation}
    \hat{p}_N(x_{0:T}\mid y_{1:T}) = \sum_{i=1}^{N} \tilde{w}_T^{(i)} \delta(x_{0:T}-X_{0:T}^{(i)}).
\end{equation}

Another important quantity is the \emph{filtering distribution}, which is the marginal of the smoothing distribution:
\begin{equation}
    p(x_T \mid y_{1:T}) = \int p(x_{0:T} \mid y_{1:T}) \,\mathrm{d}x_{0:T-1}
\end{equation}
Therefore, the filtering distribution can be approximated by integrating out the previous states
\begin{align}
    \hat{p}_N(x_T \mid y_{1:T}) &= \int \hat{p}_N(x_{0:T} \mid y_{1:T}) \,\mathrm{d}x_{0:T-1} \\
    &= \sum_{i=1}^{N} \tilde{w}_T^{(i)} \delta(x_T-X_T^{(i)}) \int \delta(x_{0:T-1}-X_{0:T-1}^{(i)}) \,\mathrm{d}x_{0:T-1} \\
    &= \sum_{i=1}^{N} \tilde{w}_T^{(i)} \delta(x_T-X_T^{(i)}).
\end{align}
Consequently, the expectation of a function $g$ under the filtering distribution, analytically given by
\begin{equation}
    I_T = E[g(X_T) \mid Y_{1:T}] = \int g(x_T)p(x_T \mid y_{1:T})\,\mathrm{d}x_T,
\end{equation}
can be approximated by
\begin{equation}
    \hat{I}_T = \sum_{i=1}^{N}\tilde{w}_{T}^{(i)} g(X_T^{(i)}).
\end{equation}

The final distribution of interest is the \emph{predictive distribution}. It is obtained by propagating the filtering distribution through the transition density:
\begin{equation}
    p(x_{T+1} \mid y_{1:T}) = \int p(x_{T+1} \mid x_T)p(x_T \mid y_{1:T})\,\mathrm{d}x_T.
\end{equation}
Using the empirical approximation of the filtering distribution,we obtain the following approximation of the predictive distribution:
\begin{align}
    \hat{p}_N(x_{T+1} \mid y_{1:T}) &= \int p(x_{T+1} \mid x_T) \hat{p}_N(x_T \mid y_{1:T})\,\mathrm{d}x_T \\
    &= \sum_{i=1}^N \tilde{w}_T^{(i)} p(x_{T+1} \mid X_T^{(i)}).
\end{align}
This expression represents a mixture of transition densities, weighted by the normalized particle weights. To obtain a particle approximation, we propagate each particle by independently sampling a new state from its transition density $X_{T+1}^{(i)} \sim p(x_{T+1} \mid X_T^{(i)})$, obtaining an empirical approximation
\begin{equation}
    \hat{p}_N(x_{T+1} \mid y_{1:T}) = \sum_{i=1}^N \tilde{w}_T^{(i)} \delta(x_{T+1}-X_{T+1}^{(i)}).
\end{equation}
Consequently, the expectation of a function $g$ under the predictive distribution, analytically given by
\begin{equation}
    I_{T+1} = E[g(X_{T+1}) \mid Y_{1:T}] = \int g(x_{T+1})p(x_{T+1} \mid y_{1:T})\,\mathrm{d}x_{T+1},
\end{equation}
can be approximated by
\begin{equation}
    \hat{I}_{T+1} = \sum_{i=1}^{N}\tilde{w}_{T}^{(i)} g(X_{T+1}^{(i)}).
\end{equation}

In summary, the sequential updates of the importance weights and the particle propagation, allows us to approximate the filtering, smoothing, and predictive distributions, and their expectations, in a sequential manner. We can sum up the sequential importance sampling in Algorithm~\ref{alg:sequential_importance_sampling}.
Note that the expectations of functions under the filtering, smoothing, and predictive distributions are all biased but consistent estimators, as they are all self-normalized. Only the marginal likelihood estimator is unbiased.
\begin{algorithm}
\caption{Sequential importance sampling}
\label{alg:sequential_importance_sampling}
\begin{algorithmic}[1]
    \REQUIRE Observations $y_{1:T}$, number of particles $N$, proposal densities $q_0,q_t$, and model densities $\chi,p,f$.
    \ENSURE Weighted particles $\{X_{0:T}^{(i)},\tilde{w}_T^{(i)}\}_{i=1}^N$.

    \FOR{$i=1,\ldots,N$}
        \STATE Sample $X_0^{(i)} \sim q_0(x_0)$.
        \STATE Set $w_0^{(i)} = \dfrac{\chi(X_0^{(i)})}{q_0(X_0^{(i)})}$.
    \ENDFOR
    \STATE Normalize weights to obtain $\tilde{w}_0^{(i)}$.

    \FOR{$t=1,\ldots,T$}
        \FOR{$i=1,\ldots,N$}
            \STATE Propagate $X_t^{(i)} \sim q_t(\cdot\mid X_{0:t-1}^{(i)},y_{1:t})$.
            \STATE Compute $\alpha_t^{(i)} = \dfrac{f(y_t\mid X_t^{(i)})p(X_t^{(i)}\mid X_{t-1}^{(i)})}{q_t(X_t^{(i)}\mid X_{0:t-1}^{(i)},y_{1:t})}$.
            \STATE Update $w_t^{(i)} = w_{t-1}^{(i)}\alpha_t^{(i)}$.
        \ENDFOR
        \STATE Normalize weights to obtain $\tilde{w}_t^{(i)}$.
    \ENDFOR
\end{algorithmic}
\end{algorithm}

The main limitation of SIS is \emph{weight degeneracy}. As the number of observations increases, the importance weights tend to become increasingly unequal, as only a small fraction of particle trajectories can explain the observations. This motivates the idea of resampling.

\section{Sequential importance sampling with resampling}
\label{sec:sequential_importance_sampling_with_resampling}
Sequential importance sampling with resampling (SISR), also known as particle filtering, is an extension of SIS that incorporates a resampling step to mitigate the problem of weight degeneracy. 
Resampling is performed based on the importance weights, when the effective sample size (ESS) falls below a certain threshold, but for our thesis we limit ourselves to resampling at every time step.

Mathematically, resampling replaces the empirical distribution
\begin{equation}
    \hat{p}_N(x_{0:t}\mid y_{1:t}) = \sum_{i=1}^{N} \tilde{w}_t^{(i)} \delta(x_{0:t}-X_{0:t}^{(i)})
\end{equation}
with a new empirical distribution
\begin{equation}
    \hat{p}_N^{\mathrm{res}}(x_{0:t}\mid y_{1:t}) = \frac{1}{N} \sum_{i=1}^{N} N^{(i)} \delta(x_{0:t}-X_{0:t}^{(i)}),
\end{equation}
where $N^{(i)}$ is the number of offspring of particle $X_{0:t}^{(i)}$ after resampling, such that $\sum_{i=1}^{N} N^{(i)} = N$. Common resampling algorithms include multinomial, stratified, residual, and systematic resampling. The latter is used in this thesis.

The validity of resampling can be established by considering its effect on estimates of expectations with respect to the weighted empirical distribution. Let $g$ be a bounded measurable function. Before resampling, its expectation under the weighted empirical distribution is
\begin{equation}
    \hat{I}_N = \sum_{i=1}^{N} \tilde{w}_t^{(i)} g(X_{0:t}^{(i)}).
\end{equation}
After resampling, the corresponding estimate becomes
\begin{equation}
    \hat{I}_N^{\mathrm{res}} = \frac{1}{N} \sum_{i=1}^{N} N^{(i)}g(X_{0:t}^{(i)}).
\end{equation}
Although $\hat{I}_N$ and $\hat{I}_N^{\mathrm{res}}$ generally differ for a particular resampling outcome, an unbiased resampling scheme satisfies
\begin{equation}
    E\left[N^{(i)}\mid \tilde{w}_t^{(1:N)}\right] = N\tilde{w}_t^{(i)}, \qquad i=1,\ldots,N,
\end{equation}
where the expectation is taken with respect to the randomness introduced by resampling, conditional on the original particles weights. 

It follows that resampling preserves expectations with respect to the weighted empirical distribution, in the sense that the resampled estimate is conditionally unbiased.

\begin{proposition}
    \label{prop:resampling_unbiased}
    Let $g$ be a bounded measurable function and $\hat{I}_N$ and $\hat{I}_N^{\mathrm{res}}$ be the corresponding estimators before and after resampling. For an unbiased resampling scheme satisfying
    \begin{equation}
        E\left[N^{(i)} \mid \tilde{w}_t^{(1:N)}\right] = N\tilde{w}_t^{(i)}, \qquad i=1,\ldots,N,
    \end{equation}
    the resampled estimator satisfies
    \begin{equation}
        E\left[\hat{I}_N^{\mathrm{res}} \mid \tilde{w}_t^{(1:N)}\right] = \hat{I}_N.
    \end{equation}
\end{proposition}
\begin{proof}
    \begin{align}
        E\left[\hat{I}_N^{\mathrm{res}} \mid \tilde{w}_t^{(1:N)}\right] &= E\left[\frac{1}{N} \sum_{i=1}^{N} N^{(i)}g(X_{0:t}^{(i)}) \,\middle|\, \tilde{w}_t^{(1:N)}\right] \\
        &= \frac{1}{N} \sum_{i=1}^{N} g(X_{0:t}^{(i)}) E\left[N^{(i)} \mid \tilde{w}_t^{(1:N)}\right] \\
        &= \frac{1}{N} \sum_{i=1}^{N} g(X_{0:t}^{(i)}) N\tilde{w}_t^{(i)} = \sum_{i=1}^{N} g(X_{0:t}^{(i)}) \tilde{w}_t^{(i)} = \hat{I}_N.
    \end{align}
\end{proof}

Resampling, creates a new set of equally weighted particles. A slight problem arises when keeping track of the entire trajectory of the latent states, as resampling mixes the trajectories of the particles. This can be solved by keeping track of the ancestry of the particles, which is done by storing the indices of the parent particles during resampling. This allows us to reconstruct the entire trajectory of each particle at any time step.
Figure~\ref{fig:resampling} illustrates the effect of resampling on the particle trajectories. The resampling step effectively prunes particles with low weights and replicates particles with high weights, to maintain a diverse set of particles that better represent the posterior distribution.
\begin{figure}
    \centering
    \begin{tikzpicture}
    \foreach \t in {0,1,2,3} {
        \foreach \j in {1,2,3,4,5} {
            \node[circle, draw, fill=gray!20] (X\t\j) at (4*\t,-1.5*\j) {$X_{\t}^{\j}$};
        }
    }
    \draw[-{Stealth[scale=1.5]}] (X01.east) -- (X11.west);
    \draw[-{Stealth[scale=1.5]}] (X02.east) -- (X12.west);
    \draw[-{Stealth[scale=1.5]}] (X03.east) -- (X13.west);
    \draw[-{Stealth[scale=1.5]}] (X04.east) -- (X14.west);
    \draw[-{Stealth[scale=1.5]}] (X04.east) -- (X15.west);
    \draw[-{Stealth[scale=1.5]}] (X12.east) -- (X21.west);
    \draw[-{Stealth[scale=1.5]}] (X12.east) -- (X22.west);
    \draw[-{Stealth[scale=1.5]}] (X12.east) -- (X23.west);
    \draw[-{Stealth[scale=1.5]}] (X14.east) -- (X24.west);
    \draw[-{Stealth[scale=1.5]}] (X15.east) -- (X25.west);
    \draw[-{Stealth[scale=1.5]}] (X21.east) -- (X31.west);
    \draw[-{Stealth[scale=1.5]}] (X21.east) -- (X32.west);
    \draw[-{Stealth[scale=1.5]}] (X24.east) -- (X33.west);
    \draw[-{Stealth[scale=1.5]}] (X24.east) -- (X34.west);
    \draw[-{Stealth[scale=1.5]}] (X25.east) -- (X35.west);

    \node[above, yshift=0.35cm] at (X01.north) {$t=0$};
    \node[above, yshift=0.35cm] at (X11.north) {$t=1$};
    \node[above, yshift=0.35cm] at (X21.north) {$t=2$};
    \node[above, yshift=0.35cm] at (X31.north) {$t=3$};
\end{tikzpicture}
    \caption{Schematic illustration of the resampling process. Particles with low weights are pruned, while particles with high weights are replicated to maintain a diverse set of particles.}
    \label{fig:resampling}
\end{figure}


\subsection{Bootstrap particle filter}
One of the simplest and most widely used particle filtering algorithms is the bootstrap particle filter, which specifies the proposal distribution as
\begin{equation}
    q_t(x_t\mid x_{0:t-1},y_{1:t}) = p(x_t\mid x_{t-1}), \qquad q_0(x_0) = \chi(x_0).
\end{equation}
where $p(x_t\mid x_{t-1})$ and $\chi(x_0)$ are the transition and initial distributions of the state space model. The importance weights are then simplified to
\begin{equation}
    w_t = w_{t-1} \frac{f(y_t\mid x_t)p(x_t\mid x_{t-1})}{q_t(x_t\mid x_{0:t-1},y_{1:t})} = w_{t-1} f(y_t\mid x_t).
\end{equation}
The entire bootstrap particle filter algorithm is presented in Algorithm~\ref{alg:bootstrap_particle_filter}.
Due to the resampling step, the resulting particles are equally weighted, so in the next time iteration, the importance weights are simply $w_t = f(y_t\mid x_t)/N$. 
Expectations with respect to the empirical distribution can be calculated either by using the weighted particles before resampling or by using the equally weighted particles after resampling. The latter approximation introduces additional Monte Carlo variability due to the randomness of the resampling step.
\begin{algorithm}
\caption{Bootstrap particle filter}
\label{alg:bootstrap_particle_filter}
\begin{algorithmic}[1]
    \REQUIRE Observations $y_{1:T}$, number of particles $N$, model densities $\chi$, $p$, and $f$.
    \ENSURE Weighted particles $\{X_{0:T}^{(i)}, \tilde{w}_T^{(i)}\}_{i=1}^N$.

    \FOR{$i=1,\ldots,N$}
        \STATE Sample $X_0^{(i)} \sim \chi(x_0)$.
        \STATE Set $w_0^{(i)} = \tilde{w}_0^{(i)} = 1/N$.
    \ENDFOR

    \FOR{$t=1,\ldots,T$}
        \FOR{$i=1,\ldots,N$}
            \STATE Propagate $X_t^{(i)} \sim p(\cdot \mid X_{t-1}^{(i)})$.
            \STATE Update $w_t^{(i)} = w_{t-1}^{(i)} f(y_t \mid X_t^{(i)})$.
        \ENDFOR
        \FOR{$i=1,\ldots,N$}
            \STATE Normalize weights $\tilde{w}_t^{(i)} = \dfrac{w_t^{(i)}}{\sum_{j=1}^N w_t^{(j)}}$.
        \ENDFOR
        \STATE Resample particles $\{X_t^{(i)}\}_{i=1}^N$ according to  $\{\tilde{w}_t^{(i)}\}_{i=1}^N$.
        \STATE Set $w_t^{(i)} = 1/N$ for all $i=1,\ldots,N$.
    \ENDFOR
\end{algorithmic}
\end{algorithm}


\section{PMCMC}

\subsection{Particle marginal Metropolis-Hastings}

\subsection{Particle Gibbs}

\section{SMC$^2$}